\chapter{Оценка величины адсорбции в двухфахной двухкомпонентной системе}

Была получена формула, позволяющая получить величину адсорбции в смеси метана и этана из экспериментальных данных о плотности и составе фаз. В отличие от предшествующих, в этой работе не были использованы приближения идеального газа и идеального раствора, а только уравнение поверхности, изотерма адсорбции Гиббса и уравнение Гиббса-Дюгема. С помощью этой формулы можно лучше понять структуру поверхности жидкой пленки, состоящей из смеси метана и этана.

Для проверки полученной формулы был использован метод молекулярной динамики. Двухфазное моделирование методом молекулярной динамики позволяет, используя ту же систему, что и в предыдущей главе, получить значения адсорбции для смеси метана и этана. Аналогично используется модель Trappe-UA. Создается длинная вычислительная ячейка с жидким слоем в центре, выводится на равновесие и происходит расчет.

Однако, в данном случае происходит расчёт других величин и зависимостей. Интерес представляет профиль плотности, поскольку по избытку одного из компонентов в фазе рассчитывается адсорбция.

Таким образом, с помощью формулы из этой главы и данных о плотностях и составах фаз, взятых из эксперимента или молекулярной динамики, можно получить значение адсорбции компонентов на границе раздела фаз и получить представление о виде профиля плотности вблизи нее.

Первая часть этой главы посвящена описанию поверхности, с которой идет работа. Вторая \--- непосредственно выводу формулы. В третьей части будут показан метод расчета адсорбции при помощи метода молекулярной динамики и полученные результаты. В четвертой будет проведено сравнение между выведенной формулой для адсорбции и более ранними, использовавшими приближения идеального газа и идеального раствора. Результаты, представленные в данной главе, опубликованы в~\cite{RJOPC285654}.

\section{Поверхность Гиббса}

Для расчета адсорбции требуется определить поверхность, для которой происходит расчет. В данной работе это поверхность Гиббса~\cite{kalikmanov_nucleation_2013}. Это некоторая плоская поверхность с нулевой толщиной:

\begin{align}
	V_g + V_l = V,
\end{align}

где $T_g$, $T_l$ и $T$ \--- объёмы газовой, жидкой фазы и всей системы соответственно.

Поверхность неподвижна, и система находится в равновесии, то есть выполняются все три условия равновесия. Первым является термическое равновесие:

\begin{align}
	T_g = T_l = T,
\end{align}

где $T_g$, $T_l$ и $T$ \--- температуры газовой, жидкой фазы и всей системы соответственно. 

Также должно соблюдаться механическое равновесие:

\begin{align}
	P_g = P_l = P,
\end{align}

где $P_g$, $P_l$ и $P$ \--- давления газовой, жидкой фазы и всей системы соответственно. 

Последним является химическое равновесие, которое показывает равенство потоков из одной фазы в другую:

\begin{align}
	\mu_g = \mu_l,
\end{align}

где $\mu_g$ и $\mu_l$ \--- химические потенциалы газовой и жидкой фазы соответственно.

Для количества частиц в таком случае будет справедливо следующее соотношение:

\begin{align}
	N_g + N_l + N^{exc} = N,
\end{align}

где $N_g$ и $N_l$ \--- химические потенциалы газовой и жидкой фазы соответственно.

где $N_g$, $N_l$ и $N$ \--- количество молекул в газовой, жидкой фазе и всей системе соответственно, $N^{exc}$  \--- избыточное число молекул, которое и необходимо определить моделированием. Через него выражается адсорбция:

\begin{align}
	\Gamma = N^{exc}/S,
\end{align}

где $\Gamma$ \--- величина адсорбции, $S$ \--- площадь межфазной границы. 

В рассматриваемой системе имеется 2 компонента, потому соотношение избытков для описанной выше поверхности Гиббса будет линейным:

\begin{align}
	\alpha N_{me}^{exc} + \beta N_{et}^{exc} = 0
\end{align}

Для каждого из компонентов можно провести границу раздела фаз так, чтобы $N^{exc}$ = 0. Однако, проведенные таким образом плоскости не совпадут. Потому были рассмотрены 2 вида поверхностей. Первый вид \--- избыток метана положителен и равен обратному избытку этана:

\begin{align}
	\label{equalads}
	N_{me}^{exc} = -N_{et}^{exc}, \Gamma_{me} = - \Gamma_{et}
\end{align}

Второй вариант поверхности \--- эквимолярная поверхность для этана:

\begin{align}
	\label{eq:noetads}
	N_{et}^{exc} = 0, \Gamma_{et} = 0
\end{align}

Выбранное уравнение поверхности является первым из уравнений, необходимых для вывода изотермы адсорбции.

\section{Вывод формулы расчета адсорбции в двухфазной двухкомпонентной системе}

Второе уравнение \--- изотерма адсорбции Гиббса:

\begin{align}
	N_{me}^{exc} d \mu_{me} + N_{et}^{exc} d \mu_{et} + S d \sigma = 0,
\end{align}

где $\mu_{me}$ и $\mu_{et}$ \--- химические потенциалы метана и этана, $\sigma$ \--- поверхностное натяжение.

Это уравнение можно переписать для линии фазового равновесия: 

\begin{align}
	N_{me}^{exc} d \mu_{me} + N_{et}^{exc} d \mu_{et} + S d \sigma = 0,
\end{align}

Это уравнение можно переписать для линии фазового равновесия: 

\begin{align}
	N_{me}^{exc} ( \frac{d \mu_{me}}{dP})_{VLE} + N_{et}^{exc} (\frac{d \mu_{et}}{dP}) + S (\frac{d \sigma}{dP}) = 0, \label{eq:isotherm}
\end{align}

где VLE \--- vapour-liquid equilibrium, линия фазового равновесия. 

Уравнение поверхности и изотерма адсорбции Гиббса образуют линейную систему относительно избытков. Однако, производные химических потенциалов по давлению остаются неизвестными. Ранее в литературе (КАКОЙ?!) ддля их определения использовались приближения идеального газа для газовой фазы и идеального раствора для жидкой фазы. Однако, в этой работе удалось этого избежать с помощью уравнения Гиббса-Дюгема:

\begin{align}
	N_{me} d \mu_{me} + N_{et} d \mu_{et} = -SdT + VdP,
\end{align}

которое при постоянной температуре примет вид:

\begin{align}
	\label{eq:GibbsDuhem}
	N_{me} d \mu_{me} + N_{et} d \mu_{et} = VdP, 
\end{align} 

где $V$ — объем фазы, $P$ — давление, $N_{me}$ и $N_{et}$ — количества вещества метана и этана в
рассматриваемой фазе. Если переписать это уравнение в эквивалентной форме для обеих фаз, можно получить линейную систему относительно искомых производных:

\begin{align}
	n_{me}^j (\frac{d \mu_{me}}{dP})_{VLE} + n_{et}^{j} (\frac{d \mu_{et}}{dP})_{VLE} = 1, j = l, g
\end{align} 

Из равенства химических потенциалов на линии равновесия следует равенство их производных по давлению. Потому в этой системе две неизвестных. Поскольку это линейная система из двух уравнений с двумя неизвестными, она имеет единственное решение, которое примет вид:

\begin{subequations}
    \label{subeq:dmudp}
    \begin{gather}
		(\frac{d \mu_{me}}{dP})_{VLE} = \frac{-n_{et}^l-n_{et}^g}{n_{me}^ln_{et}^g - n_{me}^gn_{et}^l} \label{subeq:dmudp-1} \\
		(\frac{d \mu_{et}}{dP})_{VLE} = \frac{-n_{me}^l-n_{me}^g}{n_{me}^ln_{et}^g - n_{me}^gn_{et}^l} \label{subeq:dmudp-2}
    \end{gather}
\end{subequations}

Поскольку в правой части уравнений системы присутствуют только концентрации, можно говорить о том, что искомые производные удалось получить на основе экспериментально измеряемых или получаемых с помощью молекулярной динамики параметров, не используя приближения идеального газа и раствора.

Подставив полученные занчения производных~\cref{subeq:dmudp} в уравнение изотермы~\cref{eq:isotherm} и решив полученную вместе с уравнением поверхности~\cref{equalads} систему, можно получить следующее выражение для адсорбции:

\begin{align}
\Gamma_{me} = -\Gamma_{et} = \frac{-(d\sigma / dP)_{VLE}}{(d\mu_{me} / dP)_{VLE} - (d\mu_{et} / dP)_{VLE}} \label{eq:gamma_equal}
\end{align} 

В случае уравнения~\cref{noet} адсорбция будет вычисляться по формуле:

\begin{align}
\Gamma_{me} = \frac{-(d\sigma / dP)_{VLE}}{(d\mu_{me} / dP)_{VLE}} \label{eq:gamma_noet}
\end{align}

Теперь полученные зависимости следует проверить методом молеулярной динамики.

\section{Расчет адсорбции методом молекулярной динамики}

Метод расчета адсорбции похож на метод расчета поверхностного натяжения. Рассматривается подобная система, состоящая из жидкой пленки, состоящей из метана и этана и окруженной насыщенным паром. Метан и этан так же описываются моделью TraPPE-UA, объединяющей атом углерода с прилежащими атомами водорода в единый псевдоатом. Связь C-C так же описывается гармоническим потенциалом, межмолекулярные взаимодействия \--- потенциалом Леннард-Джонса. 

Аналогично происходит выход на равновесие до момента, когда плотности станут постоянными. Это проверяется в том числе по профилю плотности \--- зависимости плотности от координаты. На рисунке приведена левая половина такого профиля, правая будет выглядеть аналогично в силу симметрии системы. На рисунке можно видеть область с низкой плотностью, соответствующую газу, и область с высокой плотностью, где расположена жидкая пленка. Следует заметить, что плотность постоянна везде, кроме области межфазной границы, что служит свидетельством равновесия в системе. Ещё можно заметить максимум на профиле метана, который показывает наличие адсорбции в системе.

\begin{figure}[ht]
    \centerfloat{
        \includegraphics[scale=0.5]{profile_python.png}
    \caption{Часть профиля плотности смеси метана и этана при температуре 213 К и давлении 22 атм.}
    }\label{fig:profile}
\end{figure}

Аппроксимация профилей проводится с помощью следующей формулы:

\begin{align}
\rho(x) = \rho_g + \frac{\rho_l-\rho_g}{2}(tanh(\frac{x-x_1}{D})-tanh(\frac{x-x_2}{D}))
\end{align}

где $x_1$, $x_2$ — координаты эквимолярных поверхностей для выбранного компонента, D — характерная толщина межфазной поверхности, $\rho_g$ и $\rho_l$ — плотности компонента в паре и в жидкости соответственно. Следует отметить, что есть и другие способы аппроксимировать межфазную поверхность, например, через функцию ошибок:

\begin{align}
\rho(x) = \rho_g + \frac{\rho_l-\rho_g}{2}(erf(\frac{x-x_1}{D})-erf(\frac{x-x_2}{D}))
\end{align}

Однако, аппроксимация суммой гиперболических тангенсов лучше описывает межфазную поверхность.

Таким образом, при помощи аппроксимации определяются плотности жидкой и газовой фаз, а также положения и толщина межфазных поверхностей. Следует заметить, что для метана и этана получается 2 разных занчения толщины межфазной границы и 4 положения поверхностей \--- по 2 для метана и этана. Однако, для метана на толщину вляиет пик, соответствующий адсорбции, потому значение толщины межфазной границы для этана является более релевантным. Тем не менее, зависимость этой величины от давления ведет себя сходным образом для обоих веществ. Можно также отметить, что с увеличением давления и растет толщина межфазной границы, и уменьшается разница между плотностями фаз, что показывает размытие границы между фазами при приближении к критической точке. Такое же явление наблюдалось и для решеточного газа в работе ССЫЛКА.

Границы межфазных поверхностей $x_1$ и $x_2$, определенные из аппроксимации, тоже различны для метана и этана. Поверхности, рассчитанные для метана, находятся дальше от жидкости по сравнению с этаном. Поскольку ширина межфазной границы для метана тоже меньше, следовательно, исследуемая жидкая пленка окружена слоем с высоким содержанием метана.

На профиле плотности для метана также можно увидеть максимум, который превышает значения плотности как в паре, так и в жидкости, что говорит об адсорбции метана на межфазной границе.   

Количество адсорбированного компонента можно рассчитать как разность количества вещества, фактически присутствующего в расчете и количества, которое было бы в системе с "идеальным" профилем, в котором плотности жидкости и газа постоянны, а граница между ними плоская и находится в точках $x_{s1}$ и $x_{s2}$:

\begin{align}
\Gamma_i = N_i - [(x_{s2}-x_{s1}) \rho_i^l + (L_x-(x_{s2}-x_{s1})) \rho_i^g]*S_{yz} \label{eq:ads}
\end{align}

где $L_x$ — длина моделируемой ячейки, $S_{yz}$ = $L_yL_z$ – поперечное сечение моделируемой ячейки, $x_{s1}$ и $x_{s2}$ — положения плоскостей, определяющих поверхность Гиббса.

Были проведены расчеты адсорбции метана методом молекулярной динамики для давлений в диапазоне от 4 до 40~атм. Были проведены расчеты для двух видов поверхностей. В первом случае поверхность выбирается так, чтобы адсорбция этана, рассчитанная по формуле~\cref{eq:ads}, была равна нулю. Во втором - суммарная адсорбция равна нулю. Расчеты проводились для постоянной температуры, равной 213~К. Результаты представлены на рисунке~\cref{fig:adsorption}, проведено сравнение с выведенной выше формулой. 

\begin{figure}[ht]
    \centerfloat{
        \includegraphics[scale=0.75]{adsorption.png}
    \caption{Зависимость адсорбции метана на поверхности жидкой пленки, состоящей из метана и этана, от доли метана в жидкости для температуры 213 К.}
    }\label{fig:adsorption}
\end{figure}

Линии проведены с использованием уравнений~\cref{eq:gamma_equal, eq:gamma_noet}. Для вычисления производной химического потенциала по давлению была проведена аппроксимация полиномом полученной методом молекулярной динамики зависимости концентрации от давления. Это позволяет с помощью формул~\cref{subeq:dmudp-1, subeq:dmudp-2} получить значения производных для всех концентраций. Аналогично можно аппроксимировать зависимость поверхностного натяжения от давления для получения производной $(d\sigma/dP)_{VLE}$. Получилось, что производная $(d\mu_{et}/dP)_{VLE} < 0$ при том, что зависимость концентрации от давления возрастает. В этом случае для системы при данных условиях не применимы приближение идеального газа для пара и приближение идеального раствора для жидкости. Этот результат, возрастание концентрации пара от давления и убывание химического потенциала, согласуются с расчетами из уравнения состояния Брусиловского, проведенными в ССЫЛКА.


\section{Сравнение с формулами, использующими приближения идеального газа и идеального раствора}

В случае использвования приближений идеального газа и идеального раствора иначе рассчитываются формулы для производных химического потенциала по давлению. Для газа используется приближение идеального газа, и коэффициент активности считается равным единице. В результате получается следующая формула\cite{baidakov_surface_2013}:

\begin{align}
\Gamma_{me} =  -\frac{x^L_{me} (1-x^L_{me})}{RT}(\frac{d\sigma}{dx^L_{me}})_T
\end{align}

\begin{figure}[ht]
    \centerfloat{
        \includegraphics[scale=0.75]{comparison.png}
    \caption{Сравнение полученной методом молекулярной динамики изотермы адсорбции для 213~К (круги) с уравнениями, выведенными с помощью приближений идеального газа и идеального раствора(пунктирная линия) и без них для поверхности, для которой $\Gamma_{me} = - \Gamma_{et}$ (сплошная линия). }
    }\label{fig:adsorption}
\end{figure}

Для рассматриваемой смеси при исследованных условиях получается следующая картина. Формула, выведенная с использованием приближений идеального газа и идеального раствора, хорошо воспроизводит зависимость, полученную с помощью метода молекулярной динамики при низких концентрациях метана. Видно, что с повышением концентрации расхождение увеличивается. Тем временем, полученная выше формула одинаково хорошо воспроизводит значения адсорбции для различных концентраций, что особенно видно на графике зависимости для поверхности с нулевой адсорбцией этана~\cref{fig:adsorption}.




\section{Выводы к четвертой главе}

Продемонстрирован эффект обогащения межфазной поверхности метаном. Как следует из оценки положения и толщины межфазных границ для различных компонентов смеси, метан окружает исследуемую жидкую пленку. Толщина границы для метана меньше, а сама граница находится дальше в газовой фазе, чем для метана. Таким образом, концентрация метана растет быстрее, чем этана, и на поверхности исследуемой пленки оказывается слой, насыщенный метаном.

Была получена формула для расчета адсорбции в двухфазной двухкомпонентной системе. В отличие от предыдущих работ, при её выводе не были использованы приближения идеального газа и идеального раствора. Она позволяет с помощью экспериментальных данных, таких как плотности фаз и концентрации компонентов в них, получить величину адсорбции. Проведено сравнение с формулой, при выводе которой были использованы приближения идеального газа и идеального раствора. Показано, что выведенная в этой работе формула лучше описывает адсорбцию при высоких концентрациях метана. 

В результате моделирования смеси метана и этана методом молекулярной динамики получены значения адсорбции для давлений от 4 до 40 атм на изотерме, соответствующей температуре 213~К. Зависимость адсорбции от доли метана в жидкой фазе согласуется с полученной формулой. Показано сходство результатов с выведенной формулой для двух поверхностей Гиббса. 



